\chapter{System Architecture}
\label{ch:architecture}

This chapter describes the system that sits on top of the methodological substrate of
\cref{ch:methodology}. We move from the layered view (\cref{sec:layers}) through the
\mcp\ tool servers (\cref{sec:servers}), the single-agent orchestrator that is the thesis's
\emph{primary} result (\cref{sec:single}), the Portfolio-Manager multi-agent extension that
is reported as a \emph{bonus} (\cref{sec:pm}), the news corpus and its causal design
(\cref{sec:news}), and the observability layer that feeds the evaluation
(\cref{sec:observability}). Throughout, we are explicit about what is on the critical path
and what is not, because that hierarchy is itself part of the contribution: the thesis is
graded on evaluation rigour, and the architecture is sequenced to protect it.

% ---------------------------------------------------------------------
% Figure: layered architecture. MCP servers (data / analytics /
% execution) -> agent orchestrator -> evaluation harness, all sitting on
% the clock substrate.
% ---------------------------------------------------------------------
\begin{figure}[t]
  \centering
  \begin{tikzpicture}[
      >=Stealth, font=\footnotesize,
      layer/.style={draw, rounded corners, align=center, minimum height=12mm,
                    thick},
      srv/.style={draw, rounded corners, align=center, minimum width=30mm,
                  minimum height=12mm, fill=accent!8, draw=accent!55},
      band/.style={draw, rounded corners, align=center, minimum width=120mm,
                   minimum height=12mm},
    ]
    % evaluation (top)
    \node[band, fill=todored!6, draw=todored!50] (eval)
      {\textbf{Evaluation harness}\quad financial metrics $\cdot$ faithfulness $\cdot$ grounding $\cdot$ sophistication $\cdot$ regime segmentation};

    % agent layer
    \node[band, below=7mm of eval, fill=gray!8, draw=gray!55] (agent)
      {\textbf{Agent orchestrator (LangGraph)}\quad single-agent loop \;$\big|$\; PM multi-agent debate};

    % MCP server row
    \node[srv, below=7mm of agent.south west, anchor=north west, xshift=2mm] (data)
      {\mcp\ \textbf{data}\\\scriptsize prices, news\\\scriptsize \tnow-filtered wrappers};
    \node[srv, right=8mm of data] (analytics)
      {\mcp\ \textbf{analytics}\\\scriptsize indicators, regime,\\\scriptsize backtest};
    \node[srv, right=8mm of analytics] (exec)
      {\mcp\ \textbf{execution}\\\scriptsize place\_order,\\\scriptsize portfolio, PnL};

    % clock substrate (bottom)
    \node[band, below=7mm of analytics, fill=todored!10, draw=todored!60, very thick,
          minimum width=120mm] (clock)
      {\textbf{Simulation clock \tnow}\quad raw disk cache (parquet) $\rightarrow$ \code{filter\_rows} ($\le\tnow$) $\rightarrow$ tool output};

    % external data behind data server
    \node[layer, above left=6mm and -28mm of data, fill=white, draw=gray!45,
          minimum width=30mm] (ext)
      {\scriptsize external sources:\\\scriptsize yfinance, Finnhub,\\\scriptsize Alpha Vantage};

    % arrows
    \draw[->, thick] (agent) -- (eval) node[midway, right]{traces, decisions.jsonl};
    \draw[<->, thick] (agent.south) -- ++(0,-6mm) -| (data.north);
    \draw[<->, thick] (agent.south) -- ++(0,-6mm) -| (analytics.north);
    \draw[<->, thick] (agent.south) -- ++(0,-6mm) -| (exec.north);
    \draw[->, thick] (data) -- (clock);
    \draw[->, thick] (analytics) -- (clock);
    \draw[->, thick] (ext) -- (data);

    \begin{scope}[on background layer]
      \node[draw=accent!30, rounded corners, fit=(data)(analytics)(exec),
            inner sep=4mm, label={[accent!70,font=\scriptsize\bfseries]above left:MCP tool layer}] {};
    \end{scope}
  \end{tikzpicture}
  \caption[Layered system architecture]{Layered architecture. External data sources enter
  only through \mcp\ \emph{data} wrappers that enforce the \tnow\ filter; the \emph{analytics}
  and \emph{execution} servers expose indicators/regime/backtest and paper-trading
  respectively. The orchestrator composes these heterogeneous servers --- the core premise of
  the thesis topic --- and emits one structured decision record per step. The evaluation
  harness consumes those records. The clock substrate underlies every read: data is cached
  raw and filtered to $\le\tnow$ at read time, so a single cache file is reusable across
  runs at different clock positions.}
  \label{fig:architecture}
\end{figure}


\section{Layered overview}
\label{sec:layers}

\Cref{fig:architecture} shows four layers. At the base is the clock substrate of
\cref{sec:timemachine}: raw cached data filtered to $\le\tnow$ at read time. Above it sit
three \mcp\ servers --- \emph{data}, \emph{analytics}, and \emph{execution} --- each a
narrow, typed surface. Above those, the agent orchestrator composes the servers in its
reasoning loop. At the top, the evaluation harness consumes the structured decision records
the orchestrator emits. The key architectural premise, and the one that matches the thesis
topic, is that the agent reaches \emph{all} capabilities --- market data, indicators,
sentiment, portfolio state, execution --- only through \mcp\ tool calls, never through
in-process shortcuts. This is what makes the backtest/live equivalence exact: the same tool
calls, the same server code, only a different clock.

\section{The MCP tool servers}
\label{sec:servers}

\subsection{Data server}

The data server wraps external sources (yfinance for prices, Finnhub and Alpha Vantage for
news) behind tools that enforce \cref{inv:nolookahead}. Its tools expose price history,
news items, and derived metadata; each filters its output through the clock before
returning. Crucially, the wrapper --- not the external source --- owns the temporal
guarantee, so swapping or adding a data provider does not move the choke point. The decision
to wrap rather than re-implement external \mcp\ data servers reflects the topic's intent:
the agent composes heterogeneous \mcp\ servers, and the thesis's own server exposes only
what is thesis-specific.

\subsection{Analytics server}

The analytics server exposes custom technical indicators, the regime detector of
\cref{sec:regime}, and backtest utilities. Indicators are precomputed once per dataset and
read through the tool, rather than recomputed every bar, so the hot path stays light. The
indicator and sentiment formulas were adapted from the supervisor's reference teaching
code, re-implemented in our codebase with type annotations, tests, and the temporal filter
rather than imported wholesale.

\subsection{Execution server}

The execution server is an in-memory paper-trading broker: \code{place\_order},
\code{get\_portfolio}, \code{get\_pnl}. It accepts an explicit price parameter rather than
looking up a price internally, which keeps it independent of the data server (no circular
dependency) and makes the price assumption explicit and traceable in every decision record.
It enforces the portfolio constraints structurally: long-only, no leverage, non-negative
cash, and a hard per-position cap. The cap is enforced at order time by clamping the
requested quantity \emph{down} to the admissible maximum --- never up --- a direction that,
when once inverted, was the cause of a catastrophic early run and is now both fixed and
tested (\cref{sec:rigour}).

\section{Single-agent orchestrator (primary result)}
\label{sec:single}

The single-agent system is the thesis's primary, canonical result. It is a minimal
LangGraph~\cite{langgraph2024} graph implementing the loop of \cref{fig:loop}: a small,
explicit, typed state object flows through perceive, reason/act, and observe nodes, with the
clock advancing once per step. LangGraph is used here in its canonical tutorial shape ---
\code{StateGraph}, nodes, conditional edges, checkpointing --- deliberately on the simplest
possible surface, so that the framework is learned \emph{off} the evaluation critical path.
The decision to introduce LangGraph in this order, rather than starting with a complex
multi-agent graph, is recorded in the decision log as a risk-management choice.

\subsection{Sizing policy and why it matters for evaluation}

The agent operates a simple, explicit position-sizing policy, stated in its system prompt: a
base target of $10\%$ of \nav\ per position, a high-conviction target of $15\%$, and a hard
cap of $20\%$. This policy is not incidental --- it is the reference against which
faithfulness is judged (\cref{sec:faithfulness}). A judge that did not know the policy would
mistake a deliberate hold at the $10\%$ target for a reasoning failure; encoding the policy
into the judge is precisely the fix that distinguishes the v1 and v2 judges.

\subsection{Decision record schema}

After each step the orchestrator writes one JSON line to \file{decisions.jsonl}, flushed
immediately so that a crash never loses a paid-for decision. The record captures both halves
of what the reasoning metrics need: the \emph{stated} side (full chain-of-thought rationale,
intended action) and the \emph{actual} side (the executed order and fill, the regime label,
the flattened indicator values, and the raw tool outputs the agent saw). Recording the raw
tool outputs --- recent OHLCV bars, the indicator dictionary, the portfolio snapshot with
each position's percentage of \nav\ --- is what makes grounding verifiable: a claim in the
prose can be checked against the number the tool actually returned. This dual capture is the
schema-level reason the faithfulness and grounding metrics are well-defined.

\subsection{Asynchronous execution without breaking determinism}

The per-bar loop issues one LLM call per ticker. To keep wall-clock time tractable, calls
within a single date are dispatched concurrently under a bounded semaphore, but all tickers
on a date are scored against the \emph{same} start-of-date portfolio snapshot, and orders
are then executed serially in a deterministic ticker order. A non-regression test asserts
that the concurrent path produces byte-identical decisions to a serial run, so concurrency
is a pure latency optimisation with no effect on results. Transient API errors are retried
with exponential back-off; this, and a token-budget fix, were prerequisites for the long run
to complete at all (\cref{sec:rigour}).

\section{Portfolio-Manager multi-agent extension (bonus)}
\label{sec:pm}

The multi-agent Portfolio Manager (PM) is reported honestly as a \emph{bonus} on top of the
single-agent result, not as a prerequisite. This hierarchy follows directly from the
project's evaluation-first principle and from the finding of Agent Market
Arena~\cite{ama2025} that architecture matters more than backbone: the interesting question
is whether deliberation changes \emph{behaviour}, specifically whether it closes the
single-agent's exit-discipline gap. \Cref{fig:pmpipeline} shows the pipeline.

% ---------------------------------------------------------------------
% Figure: PM multi-agent pipeline (TradingAgents-style communication
% topology): analysts -> bull/bear debate -> research manager -> trader
% -> risk debate -> PM weights -> memory log.
% ---------------------------------------------------------------------
\begin{figure}[t]
  \centering
  \begin{tikzpicture}[
      >=Stealth, font=\scriptsize, node distance=8mm and 11mm,
      n/.style={draw, rounded corners, align=center, minimum height=9mm,
                minimum width=20mm, thick, fill=accent!8, draw=accent!55},
      deb/.style={n, fill=todored!8, draw=todored!55},
      pm/.style={n, fill=gray!12, draw=gray!60, minimum width=26mm},
      mem/.style={draw, rounded corners, align=center, minimum height=9mm,
                  fill=gray!6, draw=gray!50, minimum width=26mm},
    ]
    % analysts column
    \node[n] (tech) {Technical\\analyst};
    \node[n, below=of tech] (news) {News\\analyst};
    \node[n, below=of news] (risk) {Risk\\analyst};

    % debate
    \node[deb, right=14mm of tech] (bull) {Bull\\researcher};
    \node[deb, below=of bull] (bear) {Bear\\researcher};
    \node[n, right=12mm of {$(bull)!0.5!(bear)$}] (rm) {Research\\manager};

    \node[n, right=12mm of rm] (trader) {Trader};

    % risk debate
    \node[deb, right=12mm of trader, yshift=8mm] (agg) {Aggressive};
    \node[deb, below=4mm of agg] (con) {Conservative};
    \node[deb, below=4mm of con] (neu) {Neutral};

    \node[pm, below=14mm of trader] (pm) {\textbf{Portfolio Manager}\\target weights};
    \node[mem, left=12mm of pm] (memlog) {Causal memory log\\\scriptsize resolved $<\tnow$ only};

    % flows
    \draw[->] (tech) -- (bull);
    \draw[->] (news) -- (bull);
    \draw[->] (risk) -- (bear);
    \draw[->] (tech) -- (bear);
    \draw[<->, dashed, todored!70] (bull) -- (bear) node[midway, right]{debate};
    \draw[->] (bull) -- (rm);
    \draw[->] (bear) -- (rm);
    \draw[->] (rm) -- (trader);
    \draw[->] (trader) -- (agg);
    \draw[->] (trader) -- (con);
    \draw[->] (trader) -- (neu);
    \draw[->] (agg.east) -- ++(5mm,0) |- (pm.east);
    \draw[->] (con.east) -- ++(5mm,0) |- (pm.east);
    \draw[->] (neu.east) -- ++(5mm,0) |- (pm.east);
    \draw[->] (pm) -- (memlog) node[midway, above, font=\scriptsize]{append};
    \draw[->, dashed] (memlog) to[bend left=20] (rm)
      node[midway, left, font=\scriptsize\itshape]{past context};

    \begin{scope}[on background layer]
      \node[draw=accent!25, rounded corners, fit=(tech)(news)(risk),
            inner sep=2.5mm, label={[font=\scriptsize\bfseries,accent!70]above:analysts}]{};
      \node[draw=todored!25, rounded corners, fit=(agg)(con)(neu),
            inner sep=2.5mm, label={[font=\scriptsize\bfseries,todored!70]above:risk debate}]{};
    \end{scope}
  \end{tikzpicture}
  \caption[Portfolio-Manager multi-agent pipeline]{The Portfolio-Manager (PM) multi-agent
  pipeline, adopting the \textsc{TradingAgents} communication topology re-implemented over
  our \mcp\ tool layer. Role analysts each see only their role-appropriate tool outputs (the
  news analyst, crucially, sees \emph{only} news items --- see \cref{sec:news}); a bounded
  Bull/Bear debate and a research manager synthesise a thesis; a trader proposes; a bounded
  Aggressive/Conservative/Neutral risk debate stress-tests it; and a single Portfolio
  Manager emits portfolio-level target weights, which a deterministic rebalancer turns into
  orders. A causal memory log feeds back only decisions resolved strictly before \tnow, so
  the deliberation never leaks future outcomes.}
  \label{fig:pmpipeline}
\end{figure}


\subsection{Communication topology}

The PM adopts the \textsc{TradingAgents}~\cite{tradingagents2024} communication topology:
role analysts (technical, news, risk) produce reports; a bounded Bull/Bear research debate
and a research manager synthesise an investment thesis; a trader proposes a position; a
bounded Aggressive/Conservative/Neutral risk debate stress-tests it; and a single Portfolio
Manager emits portfolio-level target weights. The cyclic debate is exactly the structure
that a hand-rolled loop handles awkwardly and a graph framework handles cleanly, which is
where LangGraph earns its place in the second phase rather than the first.

We adopt the \emph{contract} --- the role decomposition and the bounded message shape --- but
do not import third-party code: the topology is re-implemented over our \mcp\ tool layer
with compact JSON artifacts at each node instead of an unbounded chat transcript, so that the
cache, the decision schema, the rebalancer, and the evaluation plumbing remain unchanged. The
PM emits strict JSON parsed into a typed target-weight object; invalid JSON, negative weights,
unknown tickers, or implied leverage fail loudly at the parser boundary.

\subsection{Deterministic rebalancing and fail-loud errors}

Target weights are turned into orders by a \emph{deterministic} rebalancer, separating the
stochastic decision (the LLM allocation) from the mechanical execution (share counts under
the long-only, no-leverage, non-negative-cash constraints). A subtle but important
correctness property is that a PM API or parsing failure is tagged explicitly
(\code{is\_error}) and an all-cash fallback is written so the artifact is inspectable;
errored decisions are then excluded from every metric. Without this tag, a silent failure
would be indistinguishable from a legitimate all-cash allocation and would quietly pollute
the faithfulness tables.

\subsection{Causal memory}

The PM keeps an append-only memory log of past decisions and their later realised returns,
so that the deliberation can reference prior lessons in the manner of
\textsc{TradingAgents}. The log is gated by the clock: \code{get\_past\_context(t\_now)}
returns only entries whose decision date is strictly before \tnow\ \emph{and} whose outcomes
have already resolved, and the context is capped to the most recent entries to keep the
prompt bounded. This makes the memory a causal utility rather than a leakage channel: a
future outcome cannot enter a past deliberation by construction.

\section{News: a causally correct sentiment corpus}
\label{sec:news}

News sentiment is the data source where causal correctness is most fragile, and it is
treated as a \emph{feature} that is disabled by default rather than a load-bearing part of
the primary result. Two distinct risks shape the design.

\paragraph{Publication date versus event date.} A story about an earnings call (the event)
may be published before or after the call. Only the \emph{publication} timestamp
(\code{published\_at}) is observable to an agent at \tnow; filtering on the event date would
leak. The corpus and its clock filter therefore key strictly on \code{published\_at}.

\paragraph{Weight-memory contamination.} As in \cref{sec:leakage}, a real corpus does not
remove the model's latent knowledge of historical outcomes; it only makes the grounding
metric meaningful, by giving the news analyst's cited claims verifiable extracts to be
checked against. We state this limitation plainly rather than overclaim.

\paragraph{Corpus construction.} Candidate sources were evaluated on whether they expose a
true publication timestamp. Finnhub and Alpha Vantage do; a crawl-based source exposes only
crawl time and is therefore unsuitable as a primary timestamp. The corpus is built once,
cached to parquet, and filtered by \code{published\_at <= t\_now} at read time, exactly like
prices --- it is a static artifact, not a live scrape. A feature flag keeps the corpus off
by default; the primary single-agent result does not depend on it. We report the real state
honestly in \cref{ch:results}, including provider limitations encountered (a free Finnhub
plan that returned no rows for the older window, and a crawl key that failed locally).

\paragraph{Evidence discipline.} A diagnosed bug (\cref{sec:rigour}) had the news analyst
fabricating sentiment from technical indicators when no news was available. The fix has two
parts: the news analyst's prompt is restricted to \emph{only} news tool outputs (no
indicators, prices, or portfolio), and when no items exist for any ticker the LLM call is
skipped entirely in favour of a structured ``sentiment unavailable'' marker that contributes
zero weight to the consensus. Sentiment without a cited, timestamped extract is treated as
hallucination, not signal.

\section{Observability}
\label{sec:observability}

Every decision is traced to Langfuse Cloud~\cite{langfuse2024} as a structured object:
inputs seen, tool calls and their real outputs, the chain-of-thought, the resulting order,
latency, token count, and cost. We use Langfuse Cloud rather than self-hosting, since the
trace volume of a thesis is far below the free allowance and self-hosting would add
operational overhead for no benefit. OpenAI calls are traced via the SDK drop-in; PM
decisions emit a parent observation with child spans for each analyst, each debate stage, and
the final allocation, so the multi-agent structure is legible in the trace tree. The
reasoning metrics of \cref{ch:evaluation} consume the same captured inputs and outputs, and
the metric scores are attached back onto traces, which yields regime-segmented dashboards at
no extra cost. Trace export is best-effort: an export timeout never fails a local run, since
the authoritative artifact is the on-disk \file{decisions.jsonl}, not the trace.

\section{LLM backbone interface}
\label{sec:backbone}

The agents run on the OpenAI API behind a thin backbone interface, so that swapping or
adding a model --- for instance a second model for an architecture-versus-backbone
comparison in the spirit of~\cite{ama2025} --- is a localised change. In development the
backbone defaults to the cheapest capable model with a prompt-hash response cache, so that
re-runs during debugging cost nothing; the more expensive configuration is reserved for the
final run. A deterministic stub backbone exists for plumbing and CI, and any run using it is
labelled as such and is never cited as a thesis result.
